Quantum technologies offer new and promising alternatives across several areas of science. One of the most prominent fields is quantum computing, whose goal is to harness purely quantum phenomena to perform computational tasks faster and more efficiently than classical computers. In this project, I present, in a clear and detailed manner, the operation and use of an NMR-based quantum computer.

These computational tasks are expressed through quantum logic gate circuits, which must be implemented according to the physical system used as a platform. In this project, the implementation of a destructive \emph{SWAP}-test gate is achieved by approximating the system’s evolution when a radiofrequency magnetic field is applied, using computational methods to reproduce the desired operation. Its fidelity is then evaluated and compared with other similar implementations from the state of the art.

With well-characterized quantum gate circuits, it becomes possible to perform computational tasks such as prime factorization using Shor’s algorithm or, in the case of this project, the measurement of Bargmann invariants using \emph{SWAP}-test and \emph{cycle}-test circuits. These are fundamental quantities for characterizing the corresponding quantum states, allowing one to determine their equivalence or detect the presence of quantum entanglement.